\chapter{Willmore energy, multiplicity and inversion}\label{sec:willmore-multiplicity-inversion}
The multiplicity inequality is due to Li--Yau \cite{LY82}, Theorem 6, p.~286. We give a direct smooth proof through the radial identity underlying Simon's equations (1.2') and (1.5) in \cite{SimW93}, p.~285. The scope is smooth compact immersions; this section does not assert a corresponding theorem for arbitrary weak surfaces or branched maps. Write $\vec H=\Delta F$ and $\mathcal W(F)=\frac14\int|\vec H|^2d\mu$.

\begin{lemma}[The differential and curvature of an inversion]\label{lem:inversion-mean-curvature-vector}
Let $F:M^2\to\mathbb R^N\setminus\{y\}$ be a smooth immersion. Write $X=F-y$, $r=|X|$, and let $R_X=I-2r^{-2}X\otimes X$ be reflection across the hyperplane orthogonal to $X$; the notation $X^\perp$ below denotes projection onto the surface normal space. The inversion $I_y(z)=(z-y)/|z-y|^2$ satisfies
\[
 dI_y=r^{-2}R_X,\quad d\widetilde\mu=r^{-4}d\mu,\quad
 \widetilde{\vec H}=r^2R_X\left(\vec H+4\frac{X^\perp}{r^2}\right).
\]
Consequently on every measurable parameter set,
\[
 \frac14\int|\widetilde{\vec H}|^2d\widetilde\mu
 =4\int\left|\frac{\vec H}{4}+\frac{X^\perp}{r^2}\right|^2d\mu,
\]
with equality also when the integrals are infinite.
\end{lemma}
\begin{proof}
Differentiate $X/r^2$ to obtain the differential. The map $R_X$ is orthogonal, so lengths scale by $r^{-2}$ and surface area by $r^{-4}$. Pullback of the Euclidean ambient metric is $r^{-4}$ times that metric, a conformal change with $\phi=-2\log r$. Lemma \ref{lem:conformal-second-fundamental-form} gives the mean curvature in that pullback metric as $r^4(\vec H+4X^\perp/r^2)$. The inversion is an isometry from the pullback metric to the target metric; pushing forward by $dI_y$ proves the formula. Squaring, multiplying by the transformed area element and retaining the factor $1/4$ proves the integral identity.
\end{proof}

\begin{lemma}[The puncture flux at a smooth immersed sheet]\label{lem:willmore-puncture-flux}
Let $F:M^2\to\mathbb R^N$ be a smooth compact immersion, possibly with smooth boundary, and let $y\notin F(\partial M)$. Its preimage is a finite set of interior points, of cardinality $m$. For small $\varepsilon>0$, remove from each such sheet the region $|F-y|<\varepsilon$. Along the resulting inner boundary, the outward conormal $\eta_\varepsilon$ of the remaining surface satisfies
\[
 \int_{\text{inner boundary}}\frac{X^T}{r^2}\cdot\eta_\varepsilon\,ds\longrightarrow-2\pi m.
\]
Near each puncture, $X^\perp=O(r^2)$; hence $X^\perp/r^2$ is bounded and locally square integrable.
\end{lemma}
\begin{proof}
An immersion is an embedding on a neighborhood of each parameter point. The preimage of $y$ is closed in the compact domain and each of its points is isolated. An infinite preimage would have an accumulation point contradicting local injectivity, so it is finite. Compactness excludes any further image points near $y$ outside the chosen sheet neighborhoods.

Translate $y$ to zero and write a sheet as a graph $F(z)=(z,u(z))$ over its tangent plane, with $u(0)=0$, $Du(0)=0$, $u(z)=O(|z|^2)$ and $Du(z)=O(|z|)$. Its radial function is $r=|z|(1+O(|z|^2))$, and its gradient has $|\nabla_Mr|=1+O(r^2)$. The small radial level is a smooth perturbation of a circle, with length $2\pi\varepsilon+O(\varepsilon^3)$. The graph expansion also gives $|X^\perp|=O(r^2)$, since the leading linear vector lies in the tangent plane and the tangent plane changes by $O(r)$.

At the inner level, $\eta_\varepsilon=-\nabla_Mr/|\nabla_Mr|$ and $X^T=r\nabla_Mr$. Thus the integrand is $-|\nabla_Mr|/\varepsilon$, and its integral is $-2\pi+O(\varepsilon^2)$ for each sheet. Sum over the finite preimage. If $m=0$, a small ball misses the image and there is no puncture flux.
\end{proof}

\begin{theorem}[A square identity for a compact immersed surface]\label{thm:willmore-radial-square-identity}
Let $F:M^2\to\mathbb R^N$ be a smooth compact immersion with smooth boundary allowed. For $y\notin F(\partial M)$ let $m=\#F^{-1}(y)$, and let $\eta$ be the outward unit conormal on the original boundary. Then
\[
 \pi m+\int_{M\setminus F^{-1}(y)}
 \left|\frac{\vec H}{4}+\frac{(F-y)^\perp}{|F-y|^2}\right|^2d\mu
 =\frac{\mathcal W(F)}4+
 \frac12\int_{\partial M}\frac{(F-y)\cdot\eta}{|F-y|^2}\,ds.
\]
The boundary integral is zero for a closed surface. All integrals retain parameter multiplicity. This is the smooth immersed form of the identity used in \cite{SimW93}, equation (1.5).
\end{theorem}
\begin{proof}
On the punctured surface, set $V=X/r^2$. Ambient tangential divergence gives
\[
 \operatorname{div}_M V=\frac2{r^2}-2\frac{|X^T|^2}{r^4}
 =2\frac{|X^\perp|^2}{r^4}.
\]
The intrinsic divergence of $V^T$ is $\operatorname{div}_MV+\vec H\cdot V$. To see the sign, for an orthonormal tangent frame, differentiating $\langle V^\perp,e_i\rangle=0$ gives $\operatorname{div}_MV^\perp=-\vec H\cdot V$. Therefore
\[
 \operatorname{div}_M(V^T)
 =2\left|\frac{\vec H}{4}+\frac{X^\perp}{r^2}\right|^2-\frac18|\vec H|^2.
\]
Integrate over the surface with the small sheet disks removed. The original boundary contributes the displayed flux, and Lemma \ref{lem:willmore-puncture-flux} gives the limiting inner flux $-2\pi m$. The same lemma and smoothness of $\vec H$ give integrability of both terms near each puncture. Letting the radii tend to zero and rearranging proves the identity. Smoothness of the image as a set at a multiple point is not required: each parameter sheet has its own puncture.
\end{proof}

\begin{theorem}[The Li--Yau multiplicity bound for smooth immersions]\label{thm:li-yau-smooth-multiplicity}
For a smooth closed immersion $F:M^2\to\mathbb R^N$ and any $y\in\mathbb R^N$,
\[
 \mathcal W(F)\geq4\pi\,\#F^{-1}(y).
\]
In particular $\mathcal W(F)\geq4\pi$ for every nonempty domain, and $\mathcal W(F)<8\pi$ implies that $F$ is an embedding. The multiplicity inequality is \cite{LY82}, Theorem 6, p.~286.
\end{theorem}
\begin{proof}
Drop the nonnegative square integral in Theorem \ref{thm:willmore-radial-square-identity}. If the domain is nonempty, choose a point of its image to obtain the first universal lower bound. If the energy is less than $8\pi$, no two distinct parameter points can have the same image, so $F$ is injective. A continuous injective map from a compact space to a Hausdorff space is a homeomorphism onto its image. Together with the immersion hypothesis and its local smooth charts, this makes $F$ a smooth embedding. The energy threshold is strict; equality at $8\pi$ does not yield injectivity by this argument.
\end{proof}

\begin{proposition}[The energy lost at an inversion pole]\label{prop:willmore-inversion-energy-defect}
Under the closed-surface hypotheses above, inversion at $y$ gives an immersion of $M\setminus F^{-1}(y)$ with finite Willmore energy and
\[
 \mathcal W(I_y\circ F)=\mathcal W(F)-4\pi m,
 \qquad m=\#F^{-1}(y).
\]
Thus energy is unchanged if the pole misses the image; if it lies on $m$ smooth sheets, inversion subtracts $4\pi$ per sheet. In the equality case $\mathcal W(F)=4\pi m$, the punctured inverted immersion is minimal.
\end{proposition}
\begin{proof}
Lemma \ref{lem:inversion-mean-curvature-vector} identifies the inverted energy with four times the square integral in Theorem \ref{thm:willmore-radial-square-identity}. With no boundary that integral is $\mathcal W(F)/4-\pi m$, proving the formula and finiteness. If it vanishes, the transformed mean curvature has zero squared integral. It is smooth on the punctured domain and must vanish everywhere there. No assertion about a smooth compactification at infinity or a classification of that minimal surface is needed.
\end{proof}

\begin{theorem}[Equality in the universal Willmore bound]\label{thm:willmore-round-sphere-rigidity}
Let $M$ be connected, closed and nonempty, and $F:M^2\to\mathbb R^N$ a smooth immersion. Then $\mathcal W(F)=4\pi$ if and only if $F$ is a parametrization without multiplicity of a round two-sphere in an affine three-plane. This is the round-sphere equality case discussed in \cite{SimW93}, pp.~281--282; a direct proof follows.
\end{theorem}
\begin{proof}
Suppose the energy is $4\pi$. Theorem \ref{thm:li-yau-smooth-multiplicity} makes $F$ an embedding. For every $q\in M$, choose the pole $y=F(q)$ in Theorem \ref{thm:willmore-radial-square-identity}. Its square integral is zero, so for every $x\ne q$,
\[
 \frac{\vec H(x)}4+\frac{(F(x)-F(q))^{\perp_x}}{|F(x)-F(q)|^2}=0.
\]
Fix $x$ and a unit tangent vector $v$. Along a geodesic $q(t)$ with $q(0)=x$, $q'(0)=v$, the immersion expansion is
$F(q(t))=F(x)+t\,dF_x(v)+\tfrac12t^2B_x(v,v)+o(t^2)$.
Taking the limit in the preceding identity gives $B_x(v,v)=\vec H(x)/2$. Polarization yields $B=\tfrac12g\otimes\vec H$ at every point.

Codazzi now implies that $\vec H$ is parallel in the normal bundle. Indeed, in orthonormal tangent directions $e_1,e_2$ the equation $(\nabla^\perp_{e_1}B)(e_2,e_2)=(\nabla^\perp_{e_2}B)(e_1,e_2)$ gives $\nabla^\perp_{e_1}\vec H=0$, and interchanging the two directions gives the other derivative. Therefore $|\vec H|=c$ is constant. It cannot be zero: if $\vec H=\Delta F=0$ on a closed connected surface, integration of $F^\alpha\Delta F^\alpha$ makes all coordinate functions constant, contradicting immersion.

Differentiating $\langle\vec H,dF(V)\rangle=0$ gives $\langle D_U\vec H,dF(V)\rangle=-\langle\vec H,B(U,V)\rangle=-(c^2/2)g(U,V)$. Since the normal derivative vanishes, $D_U\vec H=-(c^2/2)dF(U)$. It follows that $F+2\vec H/c^2=C$ is constant and $|F-C|=2/c$. The three-dimensional subspace spanned by $F-C$ and the tangent plane is fixed along every curve: derivatives of $F-C$ are tangent, and derivatives of tangent fields have normal part $B$, which is parallel to $F-C$. Parallel transport of a frame for this subspace therefore keeps it in one fixed ambient three-plane. The image lies in the round two-sphere of radius $2/c$ in that plane. Its tangent spaces agree with those of the sphere, so it is open there by the inverse function theorem, and compactness makes it closed. Connectedness of the sphere gives the full image, with no multiplicity since $F$ is injective. Conversely a round sphere of radius $R$ has $|\vec H|=2/R$ and area $4\pi R^2$, so its energy is $4\pi$.
\end{proof}

\begin{corollary}[Multiplicity bounds in the unit sphere]\label{cor:spherical-willmore-multiplicity}
For a smooth closed immersion $F:M^2\to S^N$, $N\geq3$, and a point $y\in S^N$ with $m$ preimages,
\[
 \mathcal W_1(F)\geq4\pi m.
\]
In particular a closed minimal spherical immersion with area below $8\pi$ is embedded. This is the consequence stated in \cite{LY82}, Corollary 10, p.~286.
\end{corollary}
\begin{proof}
A compact smooth two-dimensional immersed image has zero $N$-dimensional volume for $N>2$, by covering it with finitely many smooth Lipschitz charts and the dimension estimate for Lipschitz images. Choose a pole outside the image. Stereographic projection preserves multiplicities, and Proposition \ref{prop:stereographic-willmore-identity} preserves the normalized energy. Apply Theorem \ref{thm:li-yau-smooth-multiplicity}. In the minimal case the energy is area, and the same injectivity argument gives embeddedness.
\end{proof}

\begin{example}[The boundary flux of a flat disk]\label{ex:willmore-flat-disk-boundary}
For the radius-$R$ disk in a two-plane, with pole at its center, the Willmore energy and the square term in Theorem \ref{thm:willmore-radial-square-identity} both vanish, but the boundary flux is $2\pi$. Thus the identity reads $\pi=\tfrac12(2\pi)$; the closed-surface lower bound cannot be applied to a disk.
\end{example}
\begin{proof}
The plane has $\vec H=0$ and $X^\perp=0$. At the boundary $\eta=X/R$, so $X\cdot\eta/r^2=1/R$. The boundary length is $2\pi R$, giving the stated flux. There is exactly one preimage of the pole. This verifies the sign and factor in the boundary identity independently of a curved example.
\end{proof}

\begin{corollary}[From the embedded theorem to all orientable positive-genus immersions]\label{cor:immersed-positive-genus-willmore}
Let $F:M^2\to\mathbb R^3$ be a smooth closed connected immersion of an orientable surface of genus at least one. Then $\mathcal W(F)\geq2\pi^2$. Equality requires embeddedness and, after inverse stereographic projection, a conformal image of the Clifford torus. The same statement holds for immersions in $S^3$ with energy $\mathcal W_1$. This deduction uses the external Marques--Neves theorem \ref{thm:willmore-bound}, as explained in \cite{MNWill}, p.~685.
\end{corollary}
\begin{proof}
If an immersion is not injective, Theorem \ref{thm:li-yau-smooth-multiplicity}, or Corollary \ref{cor:spherical-willmore-multiplicity} in the sphere, gives energy at least $8\pi>2\pi^2$. Otherwise compactness makes it an embedding. A Euclidean embedding becomes a spherical embedding by inverse stereographic projection, with the same energy by Proposition \ref{prop:stereographic-willmore-identity}. Apply Theorem \ref{thm:willmore-bound} and its equality classification. Conversely the Clifford torus and its conformal images have that energy by the explicit calculation and conformal invariance. The deep embedded theorem is an input to this argument, not proved by the multiplicity bound.
\end{proof}
